\section{The marked comparison with the extension complex}
\label{sec:marked-comparison}
This section identifies the topology obtained by replacing the parallel
positions with one-sided bands, and places the dual surface inside that
space. The boundary marking matters: a homotopy equivalence of a complex
to a sphere alone does not justify deleting disks at specified points.

\begin{proposition}[Marked punctures and band replacement]\label{prop:marked-punctures}

Choose pairwise disjoint small open two-disks about the $n$ parallel points
in $\overline T(M)$; these points lie in genuine open two-disk neighborhoods.
Let $U(M)$ be their complement. There is a homotopy equivalence of pairs
\[
 \bigl(\mathbb {RP}^2\setminus\textstyle\bigcup_{i=1}^n\mathring B_i,
       \coprod_i\partial B_i\bigr)
 \longrightarrow
 \bigl(U(M),\coprod_i\partial B_i'\bigr)
\]
whose restrictions to the boundary circles are homeomorphisms matching the
old-line labels. Replace each removed disk by a M\"obius band, attaching its
boundary by the same circle parametrization, and denote the resulting CW
complex by $B(M)$. Then
\[
 B(M)\simeq\#^{\,n+1}\mathbb {RP}^2,
 \qquad\chi(B(M))=1-n.
\]
This statement concerns a specified operation at regular parallel points;
it does not assume that the whole extension complex is a surface.
\end{proposition}
\begin{proof}

We give the marked argument because an unmarked homotopy equivalence to
$\mathbb {RP}^2$ alone would not justify puncturing.
Work first in the oriented complex $|P|$, and fix $g$. Put
$C=\mathcal E(M/g)$, a circle with antipodal sign action, regarded as a
face poset. Its zero-dimensional cells will be denoted $\rho$ and its
one-dimensional cells $\gamma$. For a closed edge $\gamma$ define
\[
 Q_\gamma=\{g^+,g^-\}\oplus C_{\leq\gamma},\quad
 R_\gamma=|\Pi_g^{-1}(Q_\gamma)|.
\]
For a vertex $\rho$ define
\[
 Q_\rho=\{g^+,g^-,\rho\},\qquad
 A_\rho=|\Pi_g^{-1}(Q_\rho)|.
\]
Each $|Q_\gamma|$ is the suspension of an interval, hence a closed disk;
each $|Q_\rho|$ is an interval with endpoints $g^+,g^-$. By the restricted
fibre argument of Proposition~\ref{prop:derivative-count}, every $R_\gamma$ and every $A_\rho$ is contractible.
The $R_\gamma$ cover $|P|$. Adjacent sectors meet in $A_\rho$; nonadjacent
sectors meet only in the two points $[g^+],[g^-]$. These assertions hold as
literal equalities of simplicial subcomplexes, because a chain in a circle's
face poset lies over one closed edge.

We describe the local marked neighborhoods. The star of any parallel
extension $[h]$ is a cone on the rank-two circle $\mathcal E(M/h)$, by
Sturmfels--Ziegler's Lemma~4.1, p.~37~\cite{SZ1993}, and $[h]$ is minimal in $P$. Thus its star is a genuine open
two-disk neighborhood in the entire order complex. For $h\neq g$, its image
under $\Pi_g$ is the rank-two parallel vertex $\rho$ corresponding to
$v=g\cap h$. In its link $\mathcal E(M/h)$, evaluation at the cocircuit of
$v$ has two zeros and separates the circle into the positive and negative
arcs. The two zeros form the link of $[h]$ in $A_\rho$; the two closed arcs
form its links in the two adjacent $R_\gamma$. Hence $[h]$ is locally an
interior interval point of $A_\rho$ and a boundary half-disk point of each
adjacent sector. Similarly, $[g^\pm]$ is an endpoint of each $A_\rho$ and a
boundary point of each sector. These local claims use the actual rank-two
circle, not a homotopy description of a link.

We next prove that the other marked points on a meridian occur between its
two endpoints. Fix $[h]\in A_\rho$. Choose an old vertex $w\in h$, $w\neq v$.
Such a vertex exists: otherwise every old line would meet $h$ at $v$, so the
whole arrangement would be concurrent. Choose the signed cocircuit $Y_w$
with $(Y_w)_g=+$. The localization evaluation at $Y_w$, interpreted as
$-1,0,1$ and extended linearly over every order simplex, defines a continuous
map
\[
 e_w:A_\rho\longrightarrow[-1,1].
\]
A simplex cannot have both a positive and a negative value, since its
vertices are comparable localizations. The only vertex of $A_\rho$ at
which this evaluation vanishes is $[h]$: every extension in the fiber
already passes through $v$, and an ordinary pseudoline cannot also pass
through $w$ on the same old line $h$. The two parallel extensions with
underlying line $h$ have different images $\rho,-\rho$, so only one lies
in this meridian. Thus $e_w^{-1}(0)=\{[h]\}$, while its values at $[g^+]$
and $[g^-]$ have opposite signs. Each such marked point separates the two
endpoints.

Here is the elementary topological fact we use to cut a meridian. If a
contractible finite CW complex has a point with an open interval
neighborhood, cutting out a smaller interval gives exactly two contractible
components whenever the cut point separates two specified points. To prove
it, restore the interval as a one-cell joining the cut complex. Its relative
homology and the long exact sequence show that the cut complex has two
acyclic components. Van Kampen identifies the fundamental group of the
restored complex with the free product of the groups of those components,
so both are simply connected. A simply connected acyclic CW complex is
contractible, by Hurewicz followed by Whitehead. Iterating proves the
statement for finitely many disjoint cuts. Since every marked point in
$A_\rho$ separates its endpoints, their separating order is linear: cutting
one places every other point on one of the two endpoint sides. No point
can be on a branch outside the endpoint-to-endpoint order.

Choose disjoint sufficiently small polygonal disks about all $2n$ oriented
parallel points, compatible with the link sectors just described and
invariant under $\tau$. Let $P^\circ$ be the complement of their interiors,
and let $R_\gamma^\circ,A_\rho^\circ$ be the induced cut pieces. Each
$R_\gamma^\circ$ is contractible. Indeed, in $R_\gamma$ each removed piece
is a boundary half-disk (a sector at $g^\pm$ is homeomorphic to a half-disk),
and the half-disk strongly deformation retracts to its frontier semicircle
while fixing that frontier. Thus deleting it leaves a deformation retract
of the original sector. The meridian $A_\rho^\circ$ is a disjoint union of
contractible components: endpoint cuts do not change its homotopy type,
and the intervening interval cuts were handled in the previous paragraph.
The cuts at $g^\pm$ eliminate all intersections of nonadjacent sectors and
all triple intersections. All pieces and intersections can be triangulated
as finite subpolyhedra, so their inclusions are cofibrations.

We now construct the comparison map with its boundary marking. Start with
the ordinary suspension sphere $S^2=\operatorname{susp}|C|$, whose meridians
are indexed by $\rho$ and whose sectors are indexed by $\gamma$. On the
$\rho$-meridian place one marked point for each $[h]\in A_\rho$, in the
separating order just proved, and mark the two suspension tips. Remove small
disks about these $2n$ marked points. We obtain a sphere $S^\circ$ with
$2n$ boundary circles. The cut sectors are disks and the cut meridians are
intervals in exactly the same incidence diagram as
$R_\gamma^\circ,A_\rho^\circ$.

Map each boundary circle homeomorphically to the corresponding actual
circle in $P^\circ$, matching its two half-circle pieces at a non-tip
parallel point and its sector pieces at a tip. For each cut meridian
interval, its two endpoints have now been prescribed in the corresponding
contractible component of $A_\rho^\circ$; choose a path between them in
that component. These paths give compatible maps on all sector boundaries.
Each such boundary map extends over the corresponding sector disk because
$R_\gamma^\circ$ is contractible. Gluing gives a continuous map
$S^\circ\to P^\circ$ which is a homeomorphism on every marked boundary
circle. Its restriction to every sector is a homotopy equivalence between
contractible spaces, and its restriction to each intersection component is
likewise a homotopy equivalence. Gluing finite CW homotopy equivalences
along these cofibrations proves that the whole map is a homotopy
equivalence. One may verify this gluing assertion inductively by writing
each addition of a sector as a pushout along the disjoint union of its
intersection components; such a pushout is a homotopy pushout because all
inclusions are cofibrations.

All choices can be made equivariantly. The sign involution exchanges the
tips, sends $\rho$ to its antipodal vertex, and reverses each meridian's
endpoint order. On the model sphere choose the action $(z,t)\mapsto(-z,-t)$;
place paired marked points at opposite heights on antipodal meridians.
Choose the boundary maps, meridian paths, and sector fillings for one
representative of each antipodal pair and transport them by $\tau$. No
sector or meridian is invariant under this action, so these prescriptions
have no fixed-piece compatibility obstruction. The resulting equivariant
map descends to a homotopy equivalence of quotient pairs: its lift is a
homotopy equivalence, its induced map on the double-cover deck groups is
the identity, and the covering-space exact sequence and Whitehead give
the quotient conclusion. The model quotient is exactly
$\mathbb {RP}^2$ with $n$ open disks removed, not an unspecified quotient
of a homotopy sphere.

Finally attach a M\"obius band to every boundary circle on both sides of
this pair equivalence, using the same parametrizations. The boundary maps
are homeomorphisms and the attachments are cofibrations; the same homotopy
pushout argument extends the equivalence. On the model surface this is the
real point blow-up at each of $n$ points, giving the connected sum of
$n+1$ projective planes. Each replacement lowers Euler characteristic by
one, which also proves the displayed Euler formula directly.
\end{proof}

Replacing the entire parallel polygon by a M\"obius band gives the
same space as the small-disk operation: the annulus between the small
disk and the polygon boundary is a collar of the inserted band.
We use this model in the following statement. The dual cell structure
of a surface with four-cycle vertex links is a quadrangulation.

\begin{proposition}[The selected-square comparison]\label{prop:selected-squares}
Subdivide every ordinary cube of $\overline T(M)$ by its
coordinate midplanes. Its vertices are the centers of the original cube
faces, hence ordinary extension labels; its edges are the ordinary cover
relations. Replace each parallel polygon by the M\"obius band of Proposition~\ref{prop:marked-punctures},
subdivided as described below. In this CW model of $B(M)$ the dual
quadrangulation $\mathcal X_2^\vee$ is an embedded subcomplex containing
the entire one-skeleton.

Specifically, if $\sigma$ has distinct zero coordinates $i,j$ and resolving them
with specified signs gives $\sigma_i,\sigma_j,\sigma_{ij}$, its diamond is
\[
 \sigma<\sigma_i<\sigma_{ij},\qquad
 \sigma<\sigma_j<\sigma_{ij}.
\]
The associated quadrangle is the union of these two closed order-complex
triangles. It occurs in the ordinary part of the dual surface map
exactly when the specified transition ports at $i,j$ are adjacent in the
full boundary cycle of its disk. This construction specifies the ordinary
quadrangles and their attaching maps. The exceptional part consists of
the entire subdivided M\"obius bands. Their rectangular cells are indexed
by the original flags $(v,e)$, and their four vertices are
$U_e^+,S_v^+,S_v^-,U_e^-$ in that cyclic order.

Thus there are specified comparison maps
\[
 \mathcal X_2\cong\mathcal X_2^\vee
 \xrightarrow{\ \iota\ }B(M)
 \xrightarrow{\ \beta\ }\overline T(M).
\]
Here $\iota$ is a cellular subcomplex inclusion and induces a surjection on
fundamental groups; $\beta$ is the continuous blow-down collapsing each inserted M\"obius core circle
to its parallel point. The inclusion $\iota$ is a homotopy equivalence
exactly when no ordinary position passes through three old vertices.
In that case it is an identification of the indicated CW spaces.
\end{proposition}
\begin{proof}
We first identify the subdivision precisely. Give the cube of an ordinary
localization $\sigma$ its coordinates in $[-1,1]^{\nu(\sigma)}$, one for
each zero cocircuit pair. A coordinate remains free on a face if its
zero is unresolved, and equals $+1$ or $-1$ if that zero has been
resolved. Subdivision at coordinate zero therefore gives vertices
indexed by the faces of this original cube. A small $k$-cube is indexed
by an interval of length $k$ in its face order: its least extension has
$k$ specified zeros that may be resolved, and its greatest extension
resolves all of them with the indicated signs. These coordinate
descriptions restrict identically on every common original face, so
the subdivisions glue over the whole ordinary cubical subcomplex.
In particular the edges are exactly cover relations, and the small
two-cubes are exactly incidence diamonds. Triangulating each diamond
along its least-to-greatest diagonal gives its two order simplices.

For the assertion about selection, the two displayed chains are distinct
simplices of the order complex. Their common edge is
$[\sigma,\sigma_{ij}]$, their other edges give the four sides of the
diamond, and their union is an embedded closed disk. In the disk-gluing
construction, a corner bordered by the two indicated transition intervals
is traversed by performing these two independent local resolutions in the
two possible orders. Their supports are disjoint, so the final extension
is the same $\sigma_{ij}$. Conversely an ordinary corner of this type
contains consecutive transition ports, by the definition of a corner of
the disk. This gives the criterion and the attaching map.

No interior simplex is counted twice: a triangle in such a diamond records
its unique least extension, the intermediate resolution, and the final
two-coordinate resolution. The other intermediate vertex is then uniquely
determined. Distinct diamonds therefore meet only in their common faces.
Their union, with their faces, is a simplicial subcomplex after adding the
indicated diagonals. The sign involution carries each diamond to its
opposite diamond, and its quotient remains embedded by the sign argument
in Proposition~\ref{prop:cubical-count}.

We now include the exceptional cells, so the description gives the whole
dual surface. Fix an old line $D_i$ and its corresponding parallel
polygon in $\overline T(M)$. Lemma~4.1 of Sturmfels--Ziegler identifies
its boundary with the rank-two circle $\mathcal E(M/i)$, with its order
reversed. After the subdivision just made, denote this boundary circle
by $C_i$. Its $4n_i$ vertices alternate between $S_v^\pm$ and $U_e^\pm$,
and its edges are the ordinary transitions at the flags $v\in\partial e$.
The rank-two sign involution $a_i:C_i\to C_i$ exchanges the two superscripts
and preserves this subdivision. It is a fixed-point-free half-turn.
The interval bundle
\[
 H_i=(C_i\times[-1,1])/\bigl((z,t)\sim(a_i z,-t)\bigr)
\]
is a M\"obius band. Its boundary is identified with $C_i$ by
$z\mapsto[z,1]$. The fiber over the pair $\{z,a_i z\}$ joins these two
boundary vertices. Taking one such fiber edge for every antipodal
vertex pair gives $2n_i$ edges. The complementary closed rectangles,
one per antipodal edge pair of $C_i$, give $2n_i$ two-cells. Each is
the rectangle with corners
$U_e^+,S_v^+,S_v^-,U_e^-$ for one flag $(v,e)$.

The corridor isotopy and endpoint calculation in Proposition~\ref{prop:surface-counts} identify these
fiber edges with the exceptional sweep transitions and these rectangles
with their four-corner circuits. In particular this is equality of
the attaching maps on every boundary edge, not merely a numerical
comparison. Every $C_i$ embeds in the ordinary complex by Proposition~\ref{prop:cubical-count}; the band
interiors are disjoint, and each band rectangle is a literal embedded
rectangle in $H_i$. The nonconcurrency hypothesis gives $n_i\geq2$, so
the four vertices of each rectangle are distinct. Consequently the
ordinary selected diamonds together with all the $H_i$ form an embedded
CW subcomplex of the specified $B(M)$. Propositions~\ref{prop:ordinary-links} and~\ref{prop:surface-counts} supply the complete
ordinary and exceptional corner circuits, respectively, so the cell
correspondence is bijective and respects every attaching map. It gives
the asserted cellular homeomorphism from $\mathcal X_2^\vee$ to that
subcomplex. This also proves that no exceptional cells have been omitted.

The inclusion contains the entire one-skeleton: every ordinary cover
relation occurs as a transition, and all band fiber edges occur as
exceptional transitions. Building $B(M)$ from this subcomplex requires
only the remaining ordinary cells of dimension at least two. Their
attachment induces a quotient on fundamental groups, proving that
$\iota_*$ is surjective. More explicitly, its kernel is the normal
subgroup generated by the boundary loops of the unselected incidence
diamonds: these are precisely the additional two-cells, and attaching
cells of dimension at least three does not change the fundamental group.
Paths from a fixed base point to the attaching loops change their
representatives only by conjugation, so this normal-subgroup description
has no path ambiguity. The additional $k$-cells for $k\geq3$ are all the
ordinary small cubes indexed by length-$k$ face intervals above. In total
the ordinary subdivision has
\[
 N_k=2^k\sum_{j\geq k}\binom jk c_j
     =\frac{2^k}{k!}F^{(k)}(1)
\]
$k$-cells: inside an original open $j$-cube choose its $k$ varying
coordinates and their $2^k$ half-intervals, keeping the other coordinates
zero. Therefore $B(M)$ is obtained from the embedded surface by attaching
$2F''(1)-F'(1)+b$ unselected squares and $N_k$ cells of each dimension
$k\geq3$, where $b=\sum_i n_i$. Here we used Proposition~\ref{prop:surface-counts}: the surface has
$F'(1)-b$ ordinary squares and $2b$ band squares. This is a finite
relative CW description of the specified comparison, including every
attaching map through its cube-face interval.

For the blow-down map choose an equivariant
homeomorphism $h_i:C_i\to S^1$, so that
$h_i(a_i z)=-h_i(z)$. In the characteristic disk of the parallel cell
define
\[
 \beta_i([z,t])=t\,h_i(z).
\]
This is well-defined, is the specified boundary identification at
$t=1$, and collapses precisely the core circle at $t=0$. Extend it by
the identity on the ordinary complex. The resulting map $\beta$ has
exactly the stated target and agrees with the ordinary comparison on
every shared edge. Equivalently, $\overline T(M)$ is the quotient of
$B(M)$ by its $n$ disjoint core circles. In particular $\beta_*$ is
surjective and its kernel is the normal subgroup generated by these
core loops, by van Kampen.

Suppose finally that no ordinary position has three zero vertices.
Every ordinary two-cube has exactly four generization ports, no
ordinary specialization port, and no exceptional port. The last claim
follows from Proposition~\ref{prop:surface-counts}, because a position near an old line has at most one
zero vertex. Proposition~\ref{prop:lens} makes the four ports alternate between the two zero
vertices, so all four of the subdivided square's diamonds are selected.
There are no higher-dimensional ordinary cubes. Thus the subcomplex
is the whole of $B(M)$, proving the identification in this case. If a
position has at least three zeros, Theorem~\ref{thm:main} gives a strict Euler-characteristic
difference, so $\iota$ cannot be a homotopy equivalence.

There are $2\nu$ generization ports at a position with $\nu$ zeros. When
there are no intervening specialization ports, their cyclic link has
$2\nu$ adjacent pairs. In general the relevant cycle includes the
specialization ports too; the rule is adjacency in that full cycle. Neither
count is a count of the stable vertices of the $\nu$-cube, which number
$2^\nu$. Furthermore the quadrangles live in its face-incidence subdivision:
their vertices are cube-face centers of dimensions $\nu,\nu-1,\nu-1,\nu-2$.
They are not the two-dimensional faces of the original unsplit cube.
The unmodified extension quotient also does not contain the exceptional
bands by this comparison: its blow-down collapses each entire core
circle. Thus the full cellular embedding belongs to the specified real
blow-up $B(M)$, while its ordinary selected diamonds lie in
$\overline T(M)$ itself.
\end{proof}
